\subsection{规则分类与RIPPER}
\label{sec:classification-rules}
\label{sec:classification-ripper}

在分类任务中，规则的结论是类别标签。第\ref{sec:classic-overview-tree-rules}节已介绍规则与树的表示差别；这里采用有序规则，预测由首次命中的类别及默认类决定，训练则需要同时处理条件增长、剪枝与覆盖。

逐条学习规则时，一个直接困难是：为了覆盖少量例外，不断增加条件会产生庞大的规则集；过早删去条件，又可能丢失有用的局部结构。早期的AQ系列探索了从样本归纳条件规则的路线，后来的规则学习把覆盖策略与误差控制结合起来。Fürnkranz与Widmer在1994年提出\term{增量减错剪枝}（incremental reduced error pruning, IREP），将单条规则的增长和立即剪枝嵌入逐条学习过程，以避免先生成庞大规则集再统一简化的开销\scite{furnkranz1994irep}

Cohen在1995年改进IREP的剪枝评分和停止准则，并加入整组规则的反复优化，形成RIPPER（Repeated Incremental Pruning to Produce Error Reduction）。其核心不是保证找到全局最优规则，而是在规则长度、覆盖范围和错误之间进行可承受的启发式搜索。下面分别说明增长、剪枝、覆盖和优化在这一过程中承担什么作用\scite{cohen1995ripper}

\subsubsection{模型结构}
\label{sec:classification-ripper-structure}

RIPPER输出有序规则以及一个默认类别$c_0$。第$r$条规则的前件可写为
\begin{equation}
  a_r(\bm{x})=\bigwedge_{j=1}^{m_r}b_{r,j}(\bm{x}),
  \qquad a_r(\bm{x})\Longrightarrow y=c_r,
  \label{eq:classification-rule-form}
\end{equation}
其中$b_{r,j}$是一个返回真或假的测试，$m_r$为该规则的条件数。例如，$x_1>60$与$x_2=\text{是}$可以组成一个合取，但每个测试的业务含义必须由特征定义决定。规则\term{覆盖}（cover）样本，指样本满足其前件，与样本真实标签是否等于后件无关。

若规则数为$s$，预测时寻找最小的$r\in\{1,\ldots,s\}$使$a_r(\bm{x})$成立，输出$c_r$；不存在这样的$r$时输出$c_0$。二分类中，可以为目标类学习若干规则，把另一类作为默认类。多分类时，经典RIPPER通常按类别出现频率由低到高学习，逐次区分当前类与尚待区分的其他类，把最常见类放在默认位置。这种类别排序有助于为少数类显式生成规则，但不构成少数类召回率的保证。

\subsubsection{参数学习}
\label{sec:classification-ripper-learning}

RIPPER采用\term{顺序覆盖}（sequential covering）：在尚未被前面规则覆盖的工作集上学习一条规则，加入有序列表后，再处理没有命中的记录。设当前目标类别为$c_+$，属于该类的记录称为正例，其余为负例。保留原始二类训练集$D^{(0)}$用于整体评分，另用$U$表示不断缩小的覆盖工作集。区分这两个对象很重要：从$U$移除一条被错误覆盖的负例，并不是把这个错误从模型评价中抹去。

学习单条规则时，先分层随机划分增长集与剪枝集，经典设置约为$2:1$。增长从空前件开始，空前件覆盖所有样本；每次添加一个测试，使前件更严格、覆盖范围更小。测试可直接使用连续特征的阈值，无须一律在训练前离散化。Cohen的实现对缺失特征采用保守约定：涉及该特征的测试在缺失记录上不成立，因此规则需要其他已知字段来覆盖这些记录。

增长时使用源自FOIL（First Order Inductive Learner）的信息增益评分。设加入条件前覆盖$p_0$个正例、$q_0$个负例，加入后覆盖$p_1$个正例、$q_1$个负例，其中$q$不表示全章的样本量$n$。在命题规则设定下，评分为
\begin{equation}
  g=p_1\left(
  \log_2\frac{p_1}{p_1+q_1}
  -\log_2\frac{p_0}{p_0+q_0}
  \right).
  \label{eq:classification-foil-gain}
\end{equation}
括号比较加入条件前后的正例纯度，外面的$p_1$避免只看纯度而忽略保留下来的正例数。候选条件必须保留正例并满足覆盖量等实现约束；增长通常持续到不再覆盖增长集负例，或没有可接受的改进条件。不可分的噪声样本可能使前一种停止情形无法达到，算法不能为追求绝对纯净而无限增长。

剪枝则允许把规则末端的一串条件一起删掉，也就是比较增长过程中得到的各个前缀，而非只尝试删去一个条件。若某个前缀在剪枝集覆盖$p$个正例、$q$个负例，RIPPER的基本评分是
\begin{equation}
  v=\frac{p-q}{p+q},\qquad p+q>0.
  \label{eq:classification-ripper-pruning}
\end{equation}
在分母为正时，它与覆盖样本中的正确率$p/(p+q)$给出相同排序。剪枝选择评分最好的可接受前缀，平局可偏向较短者；覆盖不到剪枝样本时，分数没有定义，需要回退或拒绝，而不能擅自记成满分。这一评分改变了原IREP对规则的偏好，尤其避免单凭覆盖大量样本而认可纯度不足的条件。

单条规则看起来合理，也可能使整组规则过于复杂。\term{最小描述长度}（minimum description length, MDL）把这种取舍表达为：用多少比特写下规则，再用多少比特说明规则预测错了哪些训练标签。记有序规则集为$\mathcal{R}$，总描述长度可概括为
\begin{equation}
  L_{\mathrm{total}}(\mathcal{R};D^{(0)})
  =L_{\mathrm{rules}}(\mathcal{R})
   +L_{\mathrm{errors}}(D^{(0)}\mid\mathcal{R}).
  \label{eq:classification-rule-description-length}
\end{equation}
第二项按完整规则集及默认类计算，因而仍包含误覆盖记录。Cohen的实验允许总长度暂时超过此前的最小值，但超过64比特就停止增加规则，再从后向前删除会增加总长度的规则。64是实现中的停止设置，不是统计定理给出的常数，也不要求每次加规则都立即缩短总长度。

算法~\ref{alg:classification-ripper}给出二类子问题的主要流程。伪代码中的增长和剪枝采用前述评分，描述长度的编码细节由具体实现确定；多分类时在外层加入类别排序与默认类约定。

\begin{algorithm}[RIPPER的覆盖学习与规则优化]
  \label{alg:classification-ripper}
  \begin{algorithmic}[1]
    \Require 二类样本$D^{(0)}$，目标类$c_+$，优化轮数$b$
    \Ensure 有序规则集$\mathcal{R}$及默认负类
    \State $\mathcal{R}\gets()$
    \Function{Extend}{$\mathcal{R},D^{(0)}$}
      \State $U\gets D^{(0)}$中未被$\mathcal{R}$覆盖的样本
      \State $L_{\min}\gets L_{\mathrm{total}}(\mathcal{R};D^{(0)})$
      \While{$U$中仍有正例}
        \State 将$U$分层随机划为增长集$G$和剪枝集$P$
        \State 从空前件增长规则$a$，按式\eqref{eq:classification-foil-gain}添加条件
        \State 在$P$上按式\eqref{eq:classification-ripper-pruning}选择$a$的前缀
        \If{无可接受规则，或$a$不覆盖$U$中的任何正例}
          \State \textbf{停止本轮覆盖学习}
        \EndIf
        \State 把$a\Longrightarrow c_+$追加到$\mathcal{R}$
        \State $L\gets L_{\mathrm{total}}(\mathcal{R};D^{(0)})$
        \State $L_{\min}\gets\min(L_{\min},L)$
        \If{$L>L_{\min}+64$}
          \State \textbf{停止本轮覆盖学习}
        \EndIf
        \State 从$U$移除被$a$覆盖的全部样本，包括正例和负例
      \EndWhile
      \State 从后向前删除能使总描述长度降低的规则
      \State \Return $\mathcal{R}$
    \EndFunction
    \State $\mathcal{R}\gets\Call{Extend}{\mathcal{R},D^{(0)}}$
    \For{$t=1,\ldots,b$}
      \For{按顺序考察$\mathcal{R}$中的每条规则}
        \State 从空前件生成替代规则，从原前件生成修订规则
        \State 剪枝候选规则时，评价完整规则集在剪枝集上的错误
        \State 用总描述长度在原规则、替代规则、修订规则之间选择
      \EndFor
      \State $\mathcal{R}\gets\Call{Extend}{\mathcal{R},D^{(0)}}$
    \EndFor
    \State \Return $\mathcal{R}$，无规则匹配时预测负类
  \end{algorithmic}
\end{algorithm}

优化阶段是在整组规则的上下文中改写单条规则。替代规则从空前件重新增长，修订规则在原前件上继续添加条件；对候选规则的剪枝不只看它自身的纯度，还看替换后完整有序规则集的错误。随后用MDL选择保留原规则还是采用候选，必要时重新进行规则删除和剩余正例的覆盖学习。原论文把一次优化称为RIPPER，重复两次称为RIPPER2；优化次数是可配置的搜索预算，不能解释为已经完成全局最优化。

\subsubsection{小结}
\label{sec:classification-ripper-summary}

RIPPER适合需要按条件讨论预测的场景，例如审阅异常交易或核查业务规则。增长提供局部区分能力，留出剪枝降低对增长集偶然细节的依赖，覆盖避免重复处理已经命中的输入，MDL与规则优化限制整体复杂度。这些机制控制过拟合风险，但不保证在任意噪声和样本量下都得到短而准确的规则。

连续特征可由阈值直接处理；多分类和类别不平衡也不能简单归为“不支持”或“能力弱”。真正需要检查的是少数类能否在增长集与剪枝集中获得足够样本、类别排序是否合适，以及使用总体错误率会不会忽略重要的漏报。规则变多、条件变长和重叠变复杂时，预测与人工审阅的成本都会上升，不能仅凭规则形式就认为模型已经满足可审计要求。

\subsection{模型对比与总结}
\label{sec:classification-trees-comparison}
\label{sec:classification-rules-comparison}

本小节先比较ID3、C4.5与CART，再介绍改变分裂检验、工程实现或叶模型的代表性扩展。表~\ref{tab:classification-tree-comparison}按同一组问题比较三种基本方法。C4.5的主要变化是补充分裂评分、输入处理和错误估计，CART则把二叉结构与一条可验证选择的剪枝路径结合起来。两者都不能保证找到泛化误差最小的树，因而不能把一种剪枝准则描述成对另一种的无条件优越替代。

\begin{table}[htbp]
  \centering
  \small
  \begin{tabularx}{\linewidth}{@{}l*{3}{>{\RaggedRight\arraybackslash}X}@{}}
    \toprule
    方法 & 分裂评分 & 基本分支形式 & 规模控制 \\
    \midrule
    ID3 & 信息增益 & 离散取值多路分裂 & 基本版本无后剪枝 \\
    C4.5 & 增益筛选与增益率 & 离散多路、连续阈值二分 & 悲观错误估计剪枝 \\
    CART & 基尼不纯度下降 & 统一二叉分裂 & 代价复杂度路径及验证选择 \\
    \bottomrule
  \end{tabularx}
  \caption{三种决策树在基本分类设定下的比较。输入支持和停止细节随实现版本变化。}
  \label{tab:classification-tree-comparison}
\end{table}

沿着C4.5的程序体系，C5.0进一步改进树与规则生成的实现，并加入提升集成、样本权重、误分类代价和特征筛选等功能\scite{rulequest2017c50}提升集成反复学习多个分类器，改变后续学习对样本的关注，再组合各分类器的输出；AdaBoost是这一序贯重加权路线的代表构造\scite{freund1997boosting}解释集成预测时，需要同时考察多个模型。样本权重表示不同训练记录的重要性，误分类代价则区分不同错误的损失，使训练目标能够反映实际决策中的代价差异。RuleQuest提供商业版本和GPL源码版本，其公开实验报告了特定数据与硬件条件下的速度及内存改进；实际收益随数据规模、输出树还是规则集等设置变化。

若关注的不只是类别变纯，还包括分裂是否有足够的统计证据，可以采用Kass于1980年提出的\term{卡方自动交互检测}（chi-square automatic interaction detection, CHAID）。它对类别特征的取值与标签建立列联表，先合并类别分布差异不显著的取值，再比较各特征与响应的关联，按经过多重比较调整的检验结果选择分裂；证据不足时停止。剩下的取值组形成多路分支，适合把分群结果读作交叉分类，但连续变量通常需要先分组，分组方式会影响结果。这里的显著性针对相应检验的关联假设，既不是因果证据，也不能直接作为全树不出错的保证\scite{kass1980chaid}

信息增益和基尼下降还面临另一种偏差：候选切分越多，越容易偶然找到一个高分切分。\term{条件推断树}（conditional inference tree, CTree）把“选哪个特征”与“在该特征上怎样切分”分开。Hothorn、Hornik与Zeileis在2006年的框架中用条件置换分布评价特征与响应的关联：先检验当前节点中所有特征均与响应无关联的全局原假设；不能拒绝时停止，否则选择检验证据最强的特征，再为它寻找二分位置。通过同一检验尺度和多重检验控制，可以缓解不同取值数目带来的选择偏差\scite{hothorn2006ctree}

CTree的统计解释有明确边界。置换检验依赖相应原假设下的可交换性，渐近近似的可靠程度又受样本量影响；若标签与单个特征都不相关、只通过交互发生联系，节点检验也可能停止。检验驱动的生长控制为是否继续分裂提供依据，但不等于对所有分布都严格控制泛化误差，或保证小样本任务一定优于经过验证的其他树。

也可以保留树的区域划分，却改变叶节点如何预测。\term{逻辑模型树}（logistic model tree, LMT）在每个叶区域内使用Logistic回归，即用特征的线性得分表示类别对数几率，再转成类别概率。这样，树负责不同区域之间的结构变化，叶模型负责区域内部随特征变化的概率。Landwehr、Hall与Frank的2005年期刊论文基于其2003年会议工作，通过LogitBoost逐步拟合逻辑模型，并在后代节点继续细化祖先模型，而不是把每个小叶中的模型都完全独立地从零估计\scite{landwehr2005lmt}

例如在二分类叶区域$\ell$内，LMT可输出
\begin{equation}
  \hat p_\ell(y=1\mid\bm{x})
  =\frac{1}{1+\exp[-(\bm{w}_\ell^\top\bm{x}+b_\ell)]},
  \label{eq:classification-lmt-leaf}
\end{equation}
其中$\bm{w}_\ell\in\mathbb{R}^d$和$b_\ell\in\mathbb{R}$是叶模型的系数。概率在单个叶区域内部平滑变化，但跨越不同叶的边界时仍可能跳变，也不自动具有良好的校准性质。模型拟合、迭代次数选择和树剪枝共同增加训练成本，解释时也要同时阅读路径条件与叶模型系数。

这些变种分别改变分裂证据、实现或叶模型。对树结构波动的进一步处理放在第\ref{sec:classification-ensemble}节；此处的比较仍针对单棵分类树。

本小节先比较树与RIPPER的条件组织方式，再介绍单特征规则、束搜索和决策表等代表性扩展。树与规则的差别主要是条件如何共享，以及训练时怎样寻找这些条件。树在一个节点上的分裂同时决定多个分支，适合复用共同的前置判断；直接规则学习一次寻找一个局部区域，不要求不同区域具有相同的判断顺序。两者都要反复搜索，树并不是“一次训练就完成”而规则才需要迭代。若若干类别区域共享长前缀，树可能更紧凑；若局部条件各不相同，扁平规则可以避免为维持树结构而复制分支。

复杂规则学习是否真的必要，可以用只看一个特征的方法作比较。\term{单特征规则}（one rule, OneR）对每个候选特征建立“取值到多数类”的映射，再选择训练错误最少的特征。这里的“one”指一个特征：具有多个取值的特征通常对应多条if--then规则。连续特征需要形成带有最小样本量约束的区间，防止每个数值都单独成组。Holte在1993年的研究发现，在所考察的常用基准数据上，单特征规则与更复杂学习器的预测表现并不总有很大差距。这为使用简单基线检验任务难度提供了实证依据，但具体结论仍受所选数据集限制\scite{holte1993oner}

\begin{example}[一个特征可以产生多条规则]
  \label{ex:classification-oner}
  构造一个含10个样本的二分类训练集。特征$A$有$a,b,c$三个取值，对应的“类1、类2”计数分别为$(3,1)$、$(1,2)$和$(0,3)$。按每组多数类预测，得到三条规则：$A=a$时判为类1，$A=b$或$A=c$时分别判为类2，训练错误数为$1+1+0=2$。另一个二值特征$B$的两组计数均为$(2,3)$，其单特征映射的错误数为$2+2=4$。OneR在这两个候选中选择$A$，得到的是一个特征上的三条取值规则。新记录出现未见取值时，仍需要预先约定多数类回退等处理方式。
\end{example}

若一个特征不够，又希望在条件搜索中保留多种可能，可以采用Clark与Niblett于1989年提出的CN2。它结合AQ的规则覆盖思想与类似ID3的纯度评价，用\term{束搜索}（beam search）寻找条件合取：每轮扩展已有候选，只保留评分最好的$b$个候选继续搜索，其中$b$为束宽。相较于只保留一个候选，束搜索能探索更多路径，但被舍弃的条件仍可能属于最优规则，因此并不保证避开局部最优。CN2还用显著性筛选限制偶然形成的纯净小区域\scite{clark1989cn2}

原始CN2学习有序规则，Clark与Boswell在1991年的改进中进一步研究了无序规则：针对类别分别生成条件，允许多个规则共同参与预测。有序规则以首次命中解决冲突；无序规则则需要组合匹配规则的类别分布或投票权重。无序表示减少了对单一顺序的依赖，却把解释负担转移到重叠和冲突处理上，因此规则学习和预测时的聚合约定需要配套选择\scite{clark1991cn2}

规则还可以组织成查找表。\term{决策表}（decision table）选择一个特征子集$S\subseteq\{1,\ldots,d\}$，把训练样本在这些特征上的取值组合作为键，保存同键样本的多数类。训练的难点主要转向选择$S$：特征太少会把不同类别混在一起，特征太多又使组合稀疏、难以匹配新样本。Kohavi在1995年研究了通过特征子集搜索学习决策表的方法，使用验证误差评价候选组合；最佳优先搜索不断扩展当前较好的候选子集，并在若干次扩展没有改善时停止，而不是保证遍历全部子集\scite{kohavi1995decisiontable}

决策表通常只保存训练中实际出现的组合，不必枚举所有理论上可能的取值。连续值需要离散化或另外约定匹配方法，未出现的组合可回退到训练多数类。在哈希索引下，查表本身可以很快，但还要读取并构造包含$|S|$个特征的键，特征搜索的训练代价也不能忽略。它适合可用少量离散字段表达的决策；若类别由复杂交互决定，表会变大，且相近却不完全相同的输入不能自动共享统计证据。

OneR限制使用的特征数，CN2扩大条件搜索范围，RIPPER交替覆盖、剪枝和优化，决策表则共享一组索引特征。这些方法是在不同位置限制模型复杂度，不是一条规则越来越细就必然越来越准确的路线。选择时应同时检查独立验证误差、规则规模、未匹配输入和冲突处理；只有这些行为都明确，规则形式才真正有助于使用和审阅。
